\section{推理模型与 RL 革命}

\subsection{R1-Zero：Alpha Zero 时刻}

这可能是整个访谈中最重要的技术发现。DeepSeek R1-Zero 是纯 RL 训练的结果——\textbf{没有}人类偏好数据，没有 SFT，没有 RLHF。只给模型问题和可验证的答案（数学题有标准答案，代码有单元测试），让 RL 自由探索。

结果令人震惊：``等一下让我检查一下''、自我纠正、分步推理等行为\textbf{自然涌现}——不是被人类教出来的，而是 RL 自己发现这些策略能提高奖励。

\begin{importantbox}{语言模型的 Alpha Zero 时刻}
这与 Alpha Zero 在围棋中的发现完全平行：
\begin{itemize}
    \item Alpha Go（2016）：用人类棋谱训练 + 自我博弈 $\to$ 超人水平
    \item Alpha Zero（2017）：\textbf{纯自我博弈}，零人类数据 $\to$ 远超 Alpha Go
    \item R1（2025）：用人类偏好数据训练 $\to$ 不错的推理
    \item R1-Zero（2025）：\textbf{纯 RL}，零人类偏好数据 $\to$ 更强的推理
\end{itemize}

模式完全一致：\textbf{移除人类先验使系统更强大}。

\textit{``几乎每一个深度学习中令人震惊的结果……都是试错学习。两{[}trial-and-error{]}的力量远大于一{[}imitation{]}。''} (引用 Andrej Karpathy)

RLHF 不仅是安全对齐工具——它还能提升数学/代码性能。但 R1-Zero 证明：如果你有可验证的域，连 RLHF 都不需要。
\end{importantbox}

\begin{figure}[H]
\centering
\includegraphics[width=0.92\textwidth]{frame-03-r1-zero.jpg}
\caption{R1-Zero 与纯 RL 涌现能力讨论\protect\footnotemark}
\end{figure}
\footnotetext{视频画面时间区间：02:40:00--02:40:12。}

\subsection{可验证域与 AGI 路径}

当前推理训练只在\textbf{可验证任务}上有效——数学证明有标准答案，代码有单元测试。但这个``可验证''的范围正在扩大：

下一个前沿：\textbf{计算机使用}和\textbf{机器人}作为``无限可验证''的沙箱。模型可以学习浏览网页、创建企业、赚钱——这些都是可验证的（银行账户余额不会说谎）。

\textit{``顿悟时刻将是模型学会如何在 Twitter 上获得几十万真实粉丝……或者通过当网红赚到一千万美元。''} (Dylan) 这不是玩笑——如果模型能在开放互联网上自主完成经济目标，那就是一种 AGI。

\subsubsection{搜索叠加 Chain-of-Thought}

O3/O1 Pro 不只是一条 Chain-of-Thought——它们并行启动\textbf{多条}推理链，选择最佳结果（可能使用 Monte Carlo Tree Search）。

Arc AGI 测试的惊人数据：
\begin{itemize}
    \item 1 个样本：$\sim$30\% 准确率
    \item 1000 个并行样本 $\to$ 80-90\% 准确率
    \item 代价：每个问题 \$5-20——比标准聊天贵 1000 倍
\end{itemize}

这展示了\textbf{推理时计算}（test-time compute）的威力：不需要更好的模型，只需要更多的推理时间和更多的并行尝试。这是一条完全不同于``训练更大模型''的 Scaling 路径。

\subsection{推理推断经济学}

推理模型的经济学与传统 LLM 根本不同——计算从训练时转移到了推理时。

\begin{knowledgebox}{KV Cache 与推理成本}
KV Cache 是 Transformer 推理的核心瓶颈——它存储所有前序 token 的压缩表示，随上下文长度\textbf{二次增长}。

关键数字：
\begin{itemize}
    \item 输出 token 比输入 token 贵 4 倍（串行处理 vs 并行处理）
    \item 推理模型生成 10,000+ 输出 token——巨大的内存压力
    \item 更少的并发用户（每用户占用更多内存）
    \item DeepSeek MLA 节省 80-90\% 注意力内存，但仍是二次的
\end{itemize}

R1 比 O1 便宜 27 倍（\$2 vs \$60/百万输出 token）。差异来源：
\begin{enumerate}
    \item OpenAI 的 75\%+ 毛利率
    \item DeepSeek 的架构创新（MLA + 底层库优化）
    \item 中国更低的运营成本
\end{enumerate}

但 DeepSeek 自己实际上无法服务这个模型——停止了新注册，速度 $<$5 tokens/sec，没有产能。第三方提供商（Together AI、Fireworks）以 DeepSeek 5-7 倍的价格提供 R1。
\end{knowledgebox}

\subsubsection{搜索叠加 Chain-of-Thought}

O3/O1 Pro 不只是一条 Chain-of-Thought——它们并行启动多条推理链，选择最佳结果（可能是 Monte Carlo Tree Search）。Arc AGI 测试：1000 个并行样本 $\to$ 80-90\% 准确率，vs 单样本的 30\%。代价：每个问题 \$5-20。

\subsubsection{成本下降曲线与 Jevons 悖论}

\begin{importantbox}{1200 倍的成本下降}
GPT-3 推理成本：\$60-70/百万 token $\to$ 3 年内降至几美分（通过 Llama 3B 等小模型），\textbf{1200 倍}的降幅。

DeepSeek 在 GPT-4 级别模型的成本趋势线上，但它是第一个达到这个价格点的。

Jevons 悖论被确认：DeepSeek 发布后 AWS H100 定价不降反\textbf{升}，H200 几乎脱销。\textit{``Scaling Laws 已死持续了一个月……然后 O1、O3、R1 出来了，现在变成'模型进步太快了'。''} (Dylan)
\end{importantbox}

\subsection{安全、审查与对齐}

\subsubsection{Anthropic 的安全立场}

Anthropic 据报道拥有比 O3 更好的推理模型但因安全考虑不发布。R1 的 Chain-of-Thought 可能令人不安：在中英文之间切换、出现乱码、然后给出正确答案。DeepSeek 降低了所有人的安全标准——类似于苏联太空计划对 NASA 的影响。

\subsubsection{审查的三个阶段}

\begin{enumerate}
    \item \textbf{预训练数据过滤}：最基础但最粗暴
    \item \textbf{后训练（RLHF）}：Llama 2 的过度对齐案例——\textit{``如何杀死一个 Python 进程？''} 因``杀死''一词被拒绝
    \item \textbf{System Prompt}：隐藏指令，Gemini 的``黑人纳粹''事件来自 Prompt 重写而非模型权重
\end{enumerate}

每个模型都有互联网偏见（略偏左）；Grok 试图纠正但基座模型仍然吸收了 r/politics。

\subsubsection{文化后门}

\textit{``对我们的国家优势来说，开源标准应该是美国的很重要。''} (Zuckerberg)

开源软件有过后门（Linux XZ 漏洞）——AI 模型可能嵌入文化/政治偏见。英式英语正在消亡因为美国 LLM 主导；拼写优化用 Z 不用 S。Character AI 的聊天机器人已经在影响情绪——如果是故意操纵呢？

\textit{``超人说服力将在超人智能之前到来。''} (引用 Sam Altman)

\subsection{本章小结}

R1-Zero 是语言模型的 Alpha Zero 时刻——纯 RL 训练涌现推理能力。推理经济学的核心是 KV Cache 内存瓶颈。成本 3 年内降 1200 倍但需求更快增长（Jevons 悖论）。安全和审查是多层次问题，每个阶段都有不同的风险特征。

